% Python Decimal (library)

Use when float's ~15-17 significant digits is not enough (e.g. compare sqrt sums, output 20+ decimals). Slower than float, so only use when needed. \texttt{prec} = total significant digits, NOT decimals; for coords up to $10^9$ (squared $10^{18}$) plus decimals, \texttt{prec = 50}--\texttt{100} is safe.

\begin{lstlisting}[language=Python]
from decimal import Decimal, getcontext, localcontext

getcontext().prec = 50        # significant digits (default 28); set once at top

# --- Construct: ALWAYS str or int, never float ---
a = Decimal('3.14')           # ok
b = Decimal(10)               # ok
c = Decimal(0.1)              # WRONG: stores binary float approximation

# --- Methods (use current context precision) ---
x.sqrt()                      # sqrt(x), x >= 0
x.exp()                       # e^x
x.ln()                        # natural log
x.log10()                     # log base 10
x ** Decimal(y)               # power
x.quantize(Decimal('0.01'))   # round to 2 decimals (context rounding mode)

# --- Example 1: compare sqrt(a)+sqrt(b) vs sqrt(c) ---
from decimal import Decimal, getcontext
getcontext().prec = 100
a, b, c = map(Decimal, input().split())
print("Yes" if a.sqrt() + b.sqrt() < c.sqrt() else "No")

# --- Example 2: local precision scope ---
with localcontext() as ctx:
    ctx.prec = 200
    ans = Decimal(2).sqrt()

# --- Example 3: fixed-decimal output ---
print(f"{ans:.20f}")          # 20 digits after point

# --- Tips ---
# - Exact string judge: read input as Decimal from start, print exact decimals.
# - float already enough if tolerance is 1e-6; don't use Decimal there.
# - round(Decimal(...)) uses ties-to-even (Python); use x.quantize(...) to
#   apply the context rounding mode.
\end{lstlisting}